% «Así lo preguntan en el examen» · Comprensión lectora (libros 3 y 4). Sale sola antes del taller de cada lección.
% --- Libro 3 ---
\examende{pu-el-1}{
\forma{«El texto trata principalmente de…»}{Se responde con el título, el comienzo y el cierre: el tema que recorre todo.}
\forma{«El título más adecuado para el texto sería…»}{El que cubre todo el texto, ni más amplio ni un detalle.}
\forma{«El propósito del autor es…»}{Contar, explicar, convencer: se anticipa en la prelectura y se confirma al leer.}}
\examende{pu-el-2}{
\forma{«La idea central del texto es…»}{Una oración completa que sirve para todos los párrafos.}
\forma{«El párrafo 3 cumple la función de…»}{Presentar, ejemplificar, objetar, concluir.}
\forma{«El orden de las ideas en el texto es…»}{Cuatro opciones con secuencias: se compara con el mapa hecho al margen.}
\forma{«¿Cuál enunciado resume mejor el párrafo 2?»}{La idea del párrafo, no su ejemplo más llamativo.}}
\examende{pu-el-3}{
\forma{«Todas las siguientes afirmaciones son ciertas EXCEPTO…»}{Tachar las que el texto dice.}
\forma{«¿Cuál de las siguientes NO es una razón dada por el autor?»}{Se ubican las razones y se marca la que sobra.}
\forma{«Es falso afirmar que…»}{Se busca la opción que contradice el texto.}}
\examende{pu-tt-1}{
\forma{«El texto es predominantemente…»}{Narrativo, descriptivo, expositivo, argumentativo o instructivo: ¿para qué fue escrito?}
\forma{«En el segundo párrafo predomina…»}{Un texto puede cambiar de tipología de un párrafo a otro.}
\forma{«La intención comunicativa es…»}{Informar, persuadir, entretener, dar instrucciones.}}
\examende{pu-tt-2}{
\forma{«Por sus características, el texto es…»}{Crónica, ensayo, reseña, editorial, carta, noticia.}
\forma{«El texto podría publicarse en…»}{Sección de opinión, revista científica, manual, suplemento literario.}
\forma{«El lector al que se dirige es…»}{Público general, especialistas, un destinatario concreto.}}
\examende{pu-tt-3}{
\forma{«El texto está organizado como…»}{Problema-solución, causa-efecto, comparación, orden cronológico.}
\forma{«¿Cuál de estas afirmaciones es una vivencia del narrador?»}{En testimonios: separar lo vivido de lo opinado.}
\forma{«El orden de los hechos fue…»}{Reconstruir la secuencia aunque el texto la cuente desordenada.}}
\examende{pu-ie-1}{
\forma{«Según el texto, ¿cuándo / dónde / cuánto…?»}{Palabra ancla y lectura fina.}
\forma{«El orden en que ocurren las etapas es…»}{Se siguen los marcadores: primero, después, por último.}
\forma{«De acuerdo con el segundo párrafo…»}{La pregunta ya te dice dónde buscar: no salgas de ahí.}}
\examende{pu-ie-2}{
\forma{«¿Cuál enunciado expresa lo mismo que…?»}{Mismo sentido con otras palabras.}
\forma{«La frase subrayada puede reemplazarse por…»}{Conservando quién, cuántos y qué tan seguro.}
\forma{«En otras palabras, el autor afirma que…»}{La correcta casi nunca repite las palabras del texto.}}
\examende{pu-ie-3}{
\forma{«Según el texto, es cierto que…»}{Solo vale lo que el texto sostiene.}
\forma{«El texto NO afirma que…»}{La opción que no aparece o que va más allá.}
\forma{«¿Cuál afirmación contradice al texto?»}{Distinto de «no lo dice»: aquí el texto dice lo contrario.}}
\examende{pu-ie-4}{
\forma{«¿Quién sostiene que…?»}{Se revisa la voz marcada al margen.}
\forma{«Según el experto citado…»}{No se confunde con lo que opina el autor.}
\forma{«¿Cuál afirmación pertenece al autor del texto?»}{La que va sin comillas ni «según».}}
\examende{pu-vo-1}{
\forma{«La palabra subrayada significa…»}{Se deduce por la frase.}
\forma{«Puede reemplazarse, sin cambiar el sentido, por…»}{Prueba del reemplazo.}
\forma{«El antónimo de la palabra, en el texto, es…»}{Primero el sentido, después su contrario.}}
\examende{pu-vo-2}{
\forma{«En el texto, la palabra X se usa con el sentido de…»}{Entre varios sentidos reales, el que encaja aquí.}
\forma{«¿En cuál oración la palabra tiene el mismo sentido que en el texto?»}{Se compara el uso, no la palabra.}
\forma{«La expresión se usa en sentido figurado porque…»}{Literal frente a figurado.}}
\examende{pu-vo-3}{
\forma{«La expresión "…" quiere decir que…»}{Nunca la lectura literal.}
\forma{«El refrán que mejor resume el párrafo es…»}{Se resume el párrafo en una idea y se busca el refrán que la dice.}
\forma{«Al decir "…", el personaje expresa…»}{Sentido de la frase más actitud de quien la dice.}}
\examende{pu-vo-4}{
\forma{«Por su formación, la palabra significa…»}{Prefijo, raíz y sufijo.}
\forma{«El prefijo de la palabra indica…»}{Negación, repetición, anterioridad, exceso.}
\forma{«Una palabra de la misma familia es…»}{Comparten la raíz, no solo unas letras.}}
\examende{pu-rf-1}{
\forma{«El pronombre "ella" se refiere a…»}{Concordancia y reemplazo.}
\forma{«La expresión "lo anterior" alude a…»}{Una idea completa, no una palabra.}
\forma{«"Su" (en "su decisión") remite a…»}{El poseedor se busca hacia atrás.}}
\examende{pu-rf-2}{
\forma{«La expresión "este fenómeno" reemplaza a…»}{Sustitución: se busca lo que nombra.}
\forma{«En la oración se omite…»}{Se repone el verbo o el sustantivo.}
\forma{«"Los primeros" y "los segundos" son…»}{Se respeta el orden en que fueron nombrados.}}
\examende{pu-cn-1}{
\forma{«El conector "sin embargo" indica…»}{La relación: contraste, causa, consecuencia.}
\forma{«Puede reemplazarse por…»}{Otro conector de la misma familia.}
\forma{«¿Qué conector completa mejor el espacio?»}{Primero la relación entre las dos ideas.}}
\examende{pu-cn-2}{
\forma{«El párrafo 4 cumple la función de…»}{Un verbo de función.}
\forma{«Las comillas en "…" se usan para…»}{Citar, ironizar, marcar un término.}
\forma{«La pregunta del primer párrafo sirve para…»}{Plantear el problema que el texto va a responder.}
\forma{«Lo que va entre paréntesis (o rayas)…»}{Aclara, ejemplifica o precisa.}}
% --- Libro 4 ---
\examende{pu-in-1}{
\forma{«Del texto se puede inferir que…»}{La conclusión con prueba; la más modesta.}
\forma{«Se deduce del segundo párrafo que…»}{La prueba está en ese párrafo.}
\forma{«¿Cuál conclusión NO se puede sacar?»}{La que no tiene apoyo.}}
\examende{pu-in-2}{
\forma{«La afirmación supone que…»}{Lo que debe ser cierto para que la frase tenga sentido.}
\forma{«El autor da por sentado que…»}{Una idea que no defiende porque cree que todos la aceptan.}
\forma{«La expresión "otra vez" (o "todavía") indica que…»}{Algo ya había pasado o sigue pasando.}}
\examende{pu-in-3}{
\forma{«La causa de… fue…»}{Directa u original: dibuja la cadena.}
\forma{«Una consecuencia de… es…»}{Lo que sigue en la cadena.}
\forma{«¿Cuál de las siguientes NO es una causa mencionada?»}{EXCEPTO sobre causas.}
\forma{«Si no hubiera ocurrido X, entonces…»}{Se sigue la cadena sin esa ficha.}}
\examende{pu-in-4}{
\forma{«La actitud del personaje es de…»}{Acciones y gestos, con la intensidad justa.}
\forma{«El personaje actúa así porque…»}{El motivo que explica todo lo que hace.}
\forma{«El gesto de… sugiere…»}{Un detalle que delata una emoción.}}
\examende{pu-ap-1}{
\forma{«¿Cuál caso ilustra mejor la idea del autor?»}{Mismo principio, otro tema.}
\forma{«Según el texto, ¿qué recomendaría el autor en esta situación?»}{Se aplica su idea al caso.}
\forma{«¿Cuál situación contradice lo que plantea el texto?»}{Un contraejemplo del principio.}}
\examende{pu-ap-2}{
\forma{«A es a B como C es a…»}{Se nombra la relación con una frase.}
\forma{«La situación descrita es comparable a…»}{La misma relación en otro escenario.}
\forma{«La comparación del texto entre X y Y busca mostrar que…»}{Qué tienen en común las dos orillas.}}
\examende{pu-fu-1}{
\forma{«El autor cita a… para…»}{Apoyar, oponerse o ilustrar.}
\forma{«La postura del autor frente a lo que afirma X es…»}{De acuerdo, en desacuerdo o con matices.}
\forma{«"Hay quienes creen…" introduce…»}{Una idea ajena que se va a discutir.}}
\examende{pu-fu-2}{
\forma{«¿Cuál fuente respalda mejor la afirmación?»}{Quién, qué interés, cómo lo sabe.}
\forma{«¿Qué debilita lo dicho por…?»}{Un interés propio o un método débil.}
\forma{«¿Cuál es la postura de cada participante?»}{A favor, en contra, intermedia.}}
\examende{pu-ag-1}{
\forma{«La tesis del texto es…»}{Lo que el autor quiere que pienses.}
\forma{«Un argumento que usa el autor es…»}{Una razón que apoya la tesis.}
\forma{«¿Cuál enunciado debilitaría (o fortalecería) el argumento?»}{El que ataca (o apoya) una de las patas.}
\forma{«El autor menciona la objeción para…»}{Responderla y reforzar su postura.}}
\examende{pu-ag-2}{
\forma{«La intención del autor es…»}{Informar, criticar, advertir, persuadir, conmover.}
\forma{«El tono del texto es…»}{Crítico, irónico, nostálgico, objetivo, preocupado.}
\forma{«La actitud del autor frente a… es de…»}{Aprobación, rechazo, duda, admiración.}}
\examende{pu-ag-3}{
\forma{«¿Cuál de las siguientes es una opinión?»}{La que valora y no se puede verificar.}
\forma{«¿Cuál es un hecho?»}{La verificable, aunque sea falsa.}
\forma{«En la oración hay un hecho y una opinión; la opinión es…»}{Se separa el dato del juicio.}}
\examende{pu-lf-1}{
\forma{«La figura literaria de la expresión es…»}{Comparación, metáfora, personificación, hipérbole.}
\forma{«Con la expresión "…" el autor quiere decir…»}{El sentido de la figura en el texto.}
\forma{«El efecto de la imagen es…»}{Intensificar, embellecer, hacer visible.}}
\examende{pu-lf-2}{
\forma{«La expresión es irónica porque…»}{Choque entre lo dicho y la situación.}
\forma{«En realidad, el autor quiere decir que…»}{Lo contrario del elogio aparente.}
\forma{«La intención de la ironía es…»}{Criticar o burlarse con humor.}}
\examende{pu-dt-1}{
\forma{«Los dos textos coinciden en…»}{El punto común más preciso.}
\forma{«A diferencia del texto 1, el texto 2…»}{La diferencia de tesis, de argumentos o de tono.}
\forma{«La relación entre los textos es de…»}{Oposición, complemento, ejemplo, ampliación.}}
\examende{pu-dt-2}{
\forma{«El autor del texto 1 respondería al texto 2 que…»}{Con la lógica del autor 1.}
\forma{«¿Con cuál afirmación estarían de acuerdo ambos autores?»}{La que se apoya en los dos textos, no la que simplemente ninguno contradice.}
\forma{«¿Qué propuesta aceptaría el autor del texto 2?»}{La coherente con su tesis y sus valores.}}

\InputIfFileExists{ejemplos-cl}{}{}
